% =====================================================================
%  project_reference.tex
%
%  A LaTeX-safe reference sheet for the quantum-walk / machine-learning
%  project.  It is NOT the thesis.  It is the thing you open in a second
%  window while writing the thesis, so that every definition, convention,
%  number and caveat is in one place, already typeset, already labelled.
%
%  HOW TO USE IT
%  -------------
%  * Compile standalone:      pdflatex notes/project_reference.tex
%    (Requires only amsmath, amssymb, booktabs, geometry, hyperref,
%     xcolor, longtable -- all in a standard TeX Live / MacTeX install.)
%
%  * Reuse a single equation in the thesis: every numbered equation has a
%    \label{eq:...}.  Copy the equation body across, or \input this file
%    and \eqref the label.
%
%  * Reuse the macros: the block marked "MACRO BLOCK" below is meant to be
%    copied verbatim into the thesis preamble.  Then notation is identical
%    in both documents and cannot drift.
%
%  * Find things fast: every reusable chunk is delimited by a comment of
%    the form
%          % >>> BLOCK: name
%          ...
%          % <<< END BLOCK: name
%    so   grep -n "BLOCK:" notes/project_reference.tex
%    prints a table of contents from the shell.
%
%  * Section \ref{sec:writeup-map} maps each stage of the write-up onto the
%    sections of this file, so you do not have to remember where anything is.
%
%  MAINTENANCE
%  -----------
%  Results tables in Sections \ref{sec:reproduced} and \ref{sec:observations}
%  are the living part of this document.  Add a row when a number is
%  produced; always record the config and seed alongside it, per CLAUDE.md
%  Sec. 9 ("Numerical claims need provenance").  The engineering log --
%  what broke, what was tried, what was decided -- lives in
%  notes/PROGRESS.md, not here.
% =====================================================================

\documentclass[11pt]{article}

\usepackage[margin=1in]{geometry}
\usepackage{amsmath}
\usepackage{amssymb}
\usepackage{booktabs}
\usepackage{longtable}
\usepackage{xcolor}
\usepackage[colorlinks=true,linkcolor=blue,urlcolor=blue]{hyperref}

% >>> BLOCK: macros
% ---------------------------------------------------------------------
% MACRO BLOCK -- copy this verbatim into the thesis preamble.
% Defining notation once, here, is what stops \Delta\theta_M and
% \Delta\theta_{max} from both appearing in the same document.
% ---------------------------------------------------------------------
\newcommand{\dtheta}{\Delta\theta}                  % discrete rotation strength
\newcommand{\dthetaM}{\Delta\theta_{M}}             % continuous rotation bound
\newcommand{\dthetac}{\Delta\theta_{c}}             % critical rotation strength
\newcommand{\Prand}{P_{r}}                          % inverse-translation probability
\newcommand{\Prandc}{P_{r,c}}                       % critical inverse-translation prob.
\newcommand{\Cop}{\hat{C}}                          % coin operator
\newcommand{\Top}{\hat{T}}                          % translation operator
\newcommand{\Tinv}{\hat{T}^{-1}}                    % inverse translation
\newcommand{\ket}[1]{\lvert #1 \rangle}
\newcommand{\bra}[1]{\langle #1 \rvert}
\newcommand{\braket}[2]{\langle #1 \vert #2 \rangle}
\newcommand{\MoI}{\mathrm{MoI}}
\newcommand{\IPR}{\mathrm{IPR}}
\newcommand{\ES}{\mathit{ES}}
\newcommand{\Tr}{\operatorname{Tr}}
\newcommand{\Pdeloc}{P(\text{deloc})}
% Editorial marker: things that are inferred or provisional, never claimed.
\newcommand{\provisional}[1]{\textcolor{red}{[provisional: #1]}}
% <<< END BLOCK: macros

\title{Quantum Walks with Classical Randomness:\\
       Project Reference Sheet}
\author{Anurag Paudel \\ \small Group of Dr.\ Chih-Chun Chien, Department of Physics, UC Merced}
\date{Last updated: 8 September 2026}

\begin{document}
\maketitle
\tableofcontents
\newpage

% =====================================================================
\section{How this maps onto the write-up}
\label{sec:writeup-map}
% =====================================================================

\begin{longtable}{p{0.36\textwidth} p{0.56\textwidth}}
\toprule
\textbf{When you are writing\dots} & \textbf{Pull from\dots} \\
\midrule
\endhead
Introduction / motivation
  & Sec.~\ref{sec:phenomenology} (what the transition looks like) and
    Sec.~\ref{sec:published} (what is already known). \\
Model and methods -- the walk
  & Sec.~\ref{sec:walk}, Sec.~\ref{sec:lattice}. Equations
    \eqref{eq:state}--\eqref{eq:evolution}. \\
Model and methods -- randomness
  & Sec.~\ref{sec:channels}. One subsection per channel, with the control
    parameter and its range. \\
Model and methods -- manual diagnostics
  & Sec.~\ref{sec:observables}. Equations
    \eqref{eq:prob}--\eqref{eq:ipr}. \\
Model and methods -- classifiers
  & Sec.~\ref{sec:classifiers}, including the two \emph{different}
    critical-point rules in Sec.~\ref{sec:critical}. \\
Data / reproducibility statement
  & Sec.~\ref{sec:data-contract}. \\
Results -- Phase 1
  & Sec.~\ref{sec:reproduced}, Table~\ref{tab:phase1}. \\
Results -- Phases 2--3
  & Sec.~\ref{sec:reproduced}, Table~\ref{tab:exponents} (fill in against
    Paper~A's Table~I). \\
Discussion -- the new science
  & Sec.~\ref{sec:observations}, and the random-translation failure mode in
    Sec.~\ref{sec:published}. \\
Figures
  & Sec.~\ref{sec:figures}. \\
\bottomrule
\end{longtable}

\noindent
\textbf{Sourcing rule.} Where this sheet states something the two papers
state, it is a citation. Where it states something measured in this project,
the table row carries the config and seed. Anything neither of those is
marked \provisional{like this} and must not be carried into the thesis
without checking.

% =====================================================================
\section{The walk}
\label{sec:walk}
% =====================================================================

% >>> BLOCK: state
A walker on a one-dimensional lattice carries two internal (``coin'' or
``spin'') states $\ket{+}$ and $\ket{-}$:
\begin{equation}
\ket{\psi} = \sum_{x}\left( \alpha_{x,+}\ket{x,+} + \alpha_{x,-}\ket{x,-} \right).
\label{eq:state}
\end{equation}
% <<< END BLOCK: state

% >>> BLOCK: coin
The coin (rotation) operator is
\begin{equation}
\Cop(\theta,\phi_1,\phi_2) =
\begin{pmatrix}
\cos\theta & e^{i\phi_1}\sin\theta \\[2pt]
e^{i\phi_2}\sin\theta & -e^{i(\phi_1+\phi_2)}\cos\theta
\end{pmatrix}.
\label{eq:coin-general}
\end{equation}
Throughout both papers $\phi_1 = \phi_2 = \pi/2$, which reduces
\eqref{eq:coin-general} to the symmetric form
\begin{equation}
\Cop(\theta) =
\begin{pmatrix}
\cos\theta & i\sin\theta \\[2pt]
i\sin\theta & \cos\theta
\end{pmatrix}.
\label{eq:coin}
\end{equation}
The default base angle is $\theta_0 = \pi/6$. Paper~A confirms that
$\dthetac$ is insensitive to $\theta_0$ provided $\theta_0$ is away from $0$
and from odd multiples of $\pi/2$.
% <<< END BLOCK: coin

% >>> BLOCK: translation
The conventional translation operator moves the two spin components in
opposite directions,
\begin{equation}
\Top\ket{\psi_x,+} = \ket{\psi_{x+1},+},
\qquad
\Top\ket{\psi_x,-} = \ket{\psi_{x-1},-},
\label{eq:translation}
\end{equation}
and its inverse, needed for the random-translation channel, is
\begin{equation}
\Tinv\ket{\psi_x,+} = \ket{\psi_{x-1},+},
\qquad
\Tinv\ket{\psi_x,-} = \ket{\psi_{x+1},-}.
\label{eq:translation-inverse}
\end{equation}
% <<< END BLOCK: translation

% >>> BLOCK: evolution
The state after $N$ time steps is
\begin{equation}
\ket{\psi(t)} = \left(\Top\,\Cop\right)^{N} \ket{\psi_0},
\label{eq:evolution}
\end{equation}
i.e.\ the coin acts \emph{first} at each step, then the translation.
% <<< END BLOCK: evolution

\paragraph{Parity of the support.}
Every step changes $x$ by exactly $\pm 1$, so after $N$ steps only sites with
the same parity as $N$ carry amplitude; the remaining half of the lattice is
identically zero. This is a property of the walk, not an artefact. It is why
raw $P(x)$ plots look like a comb, and why comparison figures are usually
drawn on the reachable sub-lattice only.

% =====================================================================
\section{Lattice conventions}
\label{sec:lattice}
% =====================================================================

% >>> BLOCK: lattice-table
The two source papers use \emph{different} lattices. Mixing them is the
single most likely source of a silent bug.

\begin{table}[h]
\centering
\begin{tabular}{lll}
\toprule
 & \textbf{Paper A} (PRE \textbf{108}, 035308) & \textbf{Paper B} (PRE \textbf{110}, 064124) \\
\midrule
Lattice       & odd, $2N+1$ sites, centred on $x=0$ & even, $2N$ sites, no $x=0$ \\
Initial state & $\tfrac{1}{\sqrt{2}}\left(\ket{0,+}+\ket{0,-}\right)$
              & $\sum_{x=\pm 1,\ \sigma=\pm}\tfrac{1}{2}\ket{x,\sigma}$ \\
Max steps     & $N$ & $N-1$ (conventional / symmetric), $N/2-1$ (split-step) \\
\bottomrule
\end{tabular}
\caption{Lattice conventions. Paper~B requires the even lattice because its
symmetric and split-step translation operators are \emph{not unitary} on a
lattice containing a central $x=0$ site.}
\label{tab:lattice}
\end{table}

\noindent
\textbf{This project follows Paper~A} -- odd lattice, centred initial state,
conventional walk -- because that is the localization-transition setting.
Lattice parity is an explicit parameter in the code, never hardcoded.
% <<< END BLOCK: lattice-table

% =====================================================================
\section{Channels of classical randomness}
\label{sec:channels}
% =====================================================================

All three channels below are \emph{temporal}: the randomness is redrawn at
every time step and is uniform in space. Paper~B additionally treats
\emph{spatially dependent} randomness, drawn once per site and frozen in
time; that is out of scope for the current phase.

% >>> BLOCK: channel-discrete
\subsection{Discrete random rotation}
\label{sec:ch-discrete}
Two coin operators, with angles
\begin{equation}
\theta_{1,2} = \theta_0 \pm \dtheta ,
\label{eq:discrete-coin}
\end{equation}
and at each time step a \emph{fair} classical coin selects one, with
$P(\theta_1)=P(\theta_2)=1/2$. The fairness is part of the definition of the
channel and is not a free parameter.

\noindent Control parameter: $\dtheta$, scanned over $0 \le \dtheta \le \theta_0$.
Beyond $\theta_0$ one of the two angles goes negative, which is a different walk.
% <<< END BLOCK: channel-discrete

% >>> BLOCK: channel-continuous
\subsection{Continuous random rotation}
\label{sec:ch-continuous}
A single coin whose angle is redrawn each step,
\begin{equation}
\theta(t) = \theta_0 + \dtheta(t),
\qquad
\dtheta(t) \sim \mathrm{Uniform}\!\left(0,\ \dthetaM\right).
\label{eq:continuous-coin}
\end{equation}
Note the support is \emph{one-sided}, $[0,\dthetaM)$, not symmetric about
$\theta_0$.

\noindent Control parameter: $\dthetaM$.
% <<< END BLOCK: channel-continuous

% >>> BLOCK: channel-translation
\subsection{Random translation}
\label{sec:ch-translation}
One fixed coin. At each time step the inverse translation
\eqref{eq:translation-inverse} is applied with probability $\Prand$ in place
of $\Top$.

\noindent Control parameter: $\Prand \in [0,0.5]$. Above $0.5$ the walk is
mirror-symmetric to $1-\Prand$ by parity, so scanning past $0.5$ duplicates
rather than extends.
% <<< END BLOCK: channel-translation

% >>> BLOCK: channel-naming
\subsection{Naming}
\label{sec:ch-naming}

\begin{table}[h]
\centering
\begin{tabular}{lll}
\toprule
\textbf{Figure label} & \textbf{Formal name} & \textbf{Code key} \\
\midrule
Pure QW          & no classical randomness      & \texttt{pure} \\
Jittered         & discrete random rotation     & \texttt{discrete\_coin} \\
Uniform Jittered & continuous random rotation   & \texttt{continuous\_coin} \\
Random Transform & random translation           & \texttt{random\_translation} \\
\bottomrule
\end{tabular}
\caption{The three naming systems in play. Figure labels are the shorthand
used in the existing plots; use the formal names in prose.}
\label{tab:naming}
\end{table}

\noindent
Line styles, kept fixed across every figure in the project so that plots from
different runs remain comparable: pure = blue solid, discrete = red dashed,
continuous = green dotted, random translation = purple dash-dotted. The
classical random walk (Sec.~\ref{sec:classical}) is drawn black dashed.
% <<< END BLOCK: channel-naming

% =====================================================================
\section{What the transition looks like}
\label{sec:phenomenology}
% =====================================================================

% >>> BLOCK: phenomenology
\begin{itemize}
  \item \textbf{Weak randomness $\Rightarrow$ delocalized.} The signature
        two-peak ballistic distribution; $\MoI \propto N^{2}$.
  \item \textbf{Strong randomness $\Rightarrow$ localized.} A single central
        peak: Gaussian-like for time-dependent randomness, exponential for
        spatially dependent randomness.
  \item \textbf{At the transition:} the discrete and continuous rotation
        channels \emph{flatten out}; the random-translation channel instead
        develops a \emph{three-peak} structure, a central peak coexisting
        with the two ballistic side peaks.
\end{itemize}

\noindent
That three-peak structure is exactly what defeated the neural-network methods
in Paper~A. It is expected behaviour of the physics, not a numerical fault.
% <<< END BLOCK: phenomenology

% =====================================================================
\section{Manual (baseline) observables}
\label{sec:observables}
% =====================================================================

These are the ground-truth comparators against which the classifiers are
judged.

% >>> BLOCK: obs-prob
\subsection{Probability distribution}
\begin{equation}
P(x) = \left|\psi_{+}(x)\right|^{2} + \left|\psi_{-}(x)\right|^{2}.
\label{eq:prob}
\end{equation}
The transition is located by eye, from the onset of flattening or of the
three-peak structure. This is the ``Human'' row of Table~\ref{tab:exponents}.
% <<< END BLOCK: obs-prob

% >>> BLOCK: obs-moi
\subsection{Moment of inertia}
\begin{equation}
\MoI(t) = \sum_{x} x^{2} P(x,t).
\label{eq:moi}
\end{equation}
A delocalized walk gives $\MoI \propto N^{2}$. The critical value is read off
as the \emph{kink} at which $\MoI$ departs from that scaling.
% <<< END BLOCK: obs-moi

% >>> BLOCK: obs-ipr
\subsection{Inverse participation ratio}
\begin{equation}
\IPR = \frac{\left(\sum_{x}\left|\psi_{+}(x)\right|^{2}\right)^{2}}
            {\sum_{x}\left|\psi_{+}(x)\right|^{4}},
\qquad \psi_{+}(x) \equiv \braket{x,+}{\psi}.
\label{eq:ipr}
\end{equation}
Paper~A writes the numerator and denominator in a bra--ket shorthand; the
form above is the same quantity written unambiguously, and is what the code
in \texttt{src/qw/observables.py} computes.
Two things about this definition are deliberate and must not be
``tidied up'':
\begin{enumerate}
  \item it is built from the \textbf{spin-up component only}, not from the
        total $P(x)$;
  \item it is in the \emph{participation ratio} orientation, so it grows as
        the distribution flattens and the critical value is its
        \textbf{maximum}.
\end{enumerate}
The IPR \textbf{cannot} distinguish a one-peak from a two-peak distribution.
It only flags flatness. Do not oversell it in the write-up.
% <<< END BLOCK: obs-ipr

% >>> BLOCK: obs-classical
\subsection{Classical baseline}
\label{sec:classical}
For contrast, the unbiased classical random walk on the same lattice has the
exact binomial distribution
\begin{equation}
P_{\mathrm{cl}}(x) = \binom{N}{k} p^{k}(1-p)^{N-k},
\qquad k = \frac{x+N}{2},
\label{eq:classical}
\end{equation}
zero on sites of the wrong parity, and for $p=1/2$ satisfies exactly
\begin{equation}
\MoI_{\mathrm{cl}} = \sum_x x^2 P_{\mathrm{cl}}(x) = N .
\label{eq:classical-moi}
\end{equation}
So the classical walk spreads \emph{diffusively}, $\sigma \sim \sqrt{N}$,
against the quantum walk's \emph{ballistic} $\sigma \sim N$. This is the
one-line statement of why the quantum walk is interesting, and equally the
statement of what strong classical randomness destroys.
% <<< END BLOCK: obs-classical

% =====================================================================
\section{Entanglement quantities (Paper B; later phase)}
\label{sec:entanglement}
% =====================================================================

% >>> BLOCK: entanglement
Not used in the current phase; recorded here so the conventions are fixed
before they are needed.

Entanglement entropy, from the reduced spin density matrix
$\rho_s = \Tr_x \rho_{\mathrm{tot}}$ with eigenvalues $\lambda_\pm$:
\begin{equation}
\ES(t) = -\Tr_{s}\!\left(\rho_s \ln \rho_s\right)
       = -\lambda_{+}\ln\lambda_{+} - \lambda_{-}\ln\lambda_{-}.
\label{eq:es}
\end{equation}
\textbf{Natural logarithm}, base $e$. Results in the literature quoted in
base 2 differ by a factor $\ln 2 \approx 0.69$; say which you are using.

The classical proxy (``overlap'') is
\begin{equation}
O(t) = \sum_{x} P_{+}(x,t)\, P_{-}(x,t),
\label{eq:overlap}
\end{equation}
which tracks the \emph{inverse} of $\ES$. Steady-state values of both are
long-time averages taken after transients have died away.
% <<< END BLOCK: entanglement

% =====================================================================
\section{Data contract}
\label{sec:data-contract}
% =====================================================================

% >>> BLOCK: data-contract
\begin{itemize}
  \item \textbf{One sample} $=$ the final-time probability distribution of one
        walk realization; shape $(2N+1,)$ on the Paper~A odd lattice.
  \item \textbf{One realization} $=$ one draw of the classical randomness,
        evolved to fixed $N$ (Monte Carlo over the classical coin).
  \item \textbf{Metadata} carried with every sample: channel, $\theta_0$, $N$,
        control value, seed.
  \item \textbf{Labels:} $0 = $ delocalized (weak randomness),
        $1 = $ localized (strong randomness). Never flipped.
  \item \textbf{Physical normalization:} $\sum_x P(x) = 1$.
  \item \textbf{ML normalization:} $P_{\max} = 1$, i.e.\ each distribution is
        divided by its own maximum, matching the image-recognition
        convention. Applied to \emph{both} training and inference data.
        Paper~A found this makes no difference to the SVM but substantially
        reduces variance for the MLP and CNN.
  \item \textbf{Sample size:} $1800$ training samples (estimates are stable
        from about $2000$; more gave no visible improvement).
        Train/test split $80/20$.
\end{itemize}

\paragraph{Training-window selection is a result, not a setting.}
The two labelled classes are drawn from ranges of the control parameter at
the two extremes. Too narrow, and the classes carry too little information
about the configurations; too wide, and transition-regime data leaks into
training and corrupts the critical-value estimate. The chosen windows must be
recorded and reported alongside any critical value derived from them.
% <<< END BLOCK: data-contract

% =====================================================================
\section{Classifiers}
\label{sec:classifiers}
% =====================================================================

All three expose the same interface: \texttt{fit}, \texttt{predict\_proba}
returning genuine two-class probabilities of shape $(n,2)$, and
\texttt{score}. Genuine probabilities are non-negotiable: the entire
critical-point method in Sec.~\ref{sec:critical} is built on them.

% >>> BLOCK: clf-svm
\subsection{SVM}
Linear \texttt{SGDClassifier}, with \textbf{modified Huber} loss --- not
hinge. Hinge produces a hard true/false decision; modified Huber yields
calibrated binary probabilities.

\noindent
\emph{Expected behaviour:} near-perfect test accuracy even at small sample
size. This is not evidence of anything except a working pipeline: at the two
extremes the problem is one-peak versus two-peak and is trivially separable.
Say so explicitly in the write-up rather than reporting it as a success.
% <<< END BLOCK: clf-svm

% >>> BLOCK: clf-mlp
\subsection{MLP}
\texttt{MLPClassifier}. Six layers in total:
\begin{equation*}
\text{input (sized to the lattice)} \rightarrow 400 \rightarrow 200
\rightarrow 100 \rightarrow 50 \rightarrow \text{output } 2 ,
\end{equation*}
with regularization $\alpha = 0.001$. Hyperparameters were explored with an
exhaustive grid search (\texttt{GridSearchCV}) over the hidden layer sizes
and $\alpha$. The layer sizes were held \emph{fixed} from $N=80$ to $N=1000$;
scaling them with lattice size gave negligible improvement.

\paragraph{Implementation caveat.} \texttt{scikit-learn}'s
\texttt{MLPClassifier} uses a \emph{single} logistic output unit for a binary
problem rather than two softmax units. This is equivalent up to
reparametrization, and \texttt{predict\_proba} still returns the required
$(n,2)$, but it is not literally the two-neuron output layer drawn in
Paper~A. Worth one sentence in the methods section.
% <<< END BLOCK: clf-mlp

% >>> BLOCK: clf-cnn
\subsection{CNN}
TensorFlow/Keras. Four one-dimensional convolutional layers, then a dense
layer with 2 neurons and a softmax.
\begin{itemize}
  \item Filter counts scale with lattice size: layer 1 has $2N+1$ filters,
        then half, quarter, eighth (rounding up).
  \item Kernel size 3, chosen because the lattice size is odd, so no padding
        is needed.
  \item Dropout after each layer: $0.5$, $0.5$, then $0.2$ on the final layer.
  \item Loss: sparse categorical cross-entropy; 10 epochs.
  \item \textbf{Random translation only:} add $L_2$ kernel regularization
        with $\lambda = 0.01$ on every layer. This was needed to control
        overfitting on the more varied distributions, and helped most at
        smaller $N$ ($80$--$500$).
\end{itemize}
Because the filter counts depend on $N$, the model must be retrained for each
lattice size. The network \emph{structure} stays fixed.
% <<< END BLOCK: clf-cnn

% >>> BLOCK: critical
\subsection{Critical-point extraction}
\label{sec:critical}
After training on the two extreme regimes, intermediate-randomness data is fed
in and the transition read off. \textbf{The rule differs by method, and the
asymmetry is deliberate:}
\begin{itemize}
  \item \textbf{SVM:} the point of \emph{maximal confusion}, where the two
        class probabilities are equal.
  \item \textbf{MLP and CNN:} these show a \emph{sudden jump} rather than a
        gradual crossover, so the estimate is the \emph{first} point at which
        $\Pdeloc$ falls below $50\%$.
\end{itemize}

\noindent
\textbf{Known artefact.} The ML methods estimate the critical value from a
small \emph{range} around the transition rather than from a single point.
This is the suspected cause of their systematically lower scaling exponents,
and should be stated as a limitation.
% <<< END BLOCK: critical

% =====================================================================
\section{Published results to reproduce}
\label{sec:published}
% =====================================================================

% >>> BLOCK: table-paperA
Power-law exponents of the critical randomness against system size,
$\dthetac \sim N^{-\alpha}$, from Paper~A, Table~I.

\begin{table}[h]
\centering
\begin{tabular}{lccc}
\toprule
\textbf{Method} & $\dthetac$ (discrete) & $\dthetac$ (continuous) & $\Prandc$ \\
\midrule
SVM    & $0.36 \pm 0.04$ & $0.25 \pm 0.02$ & $\phantom{-}0.28 \pm 0.04$ \\
MLP NN & $0.32 \pm 0.02$ & $0.25 \pm 0.01$ & $\phantom{-}0.07 \pm 0.06$ \\
CNN    & $0.39 \pm 0.07$ & $0.32 \pm 0.04$ & $-0.08 \pm 0.06$ \\
MoI    & $0.62 \pm 0.02$ & $0.36 \pm 0.02$ & $\phantom{-}0.51 \pm 0.04$ \\
Human  & $0.32 \pm 0.01$ & $0.36 \pm 0.01$ & $\phantom{-}0.69 \pm 0.03$ \\
IPR    & $0.46 \pm 0.02$ & $0.41 \pm 0.01$ & $\phantom{-}0.52 \pm 0.02$ \\
\bottomrule
\end{tabular}
\caption{Paper~A, Table~I. Scaling exponents $\alpha$ in $\dthetac \sim N^{-\alpha}$.}
\label{tab:paperA}
\end{table}
% <<< END BLOCK: table-paperA

% >>> BLOCK: failure-mode
\paragraph{The known failure mode, and why it is the opening.}
For random translation the MLP and CNN deviate systematically from the manual
methods --- the CNN exponent even goes negative. Paper~A tried splitting the
distribution into three spatial regions ($-N<x<-N/2$, $-N/2<x<N/2$,
$N/2<x<N$) and training on each separately, and found \emph{no improvement}.
That experiment does not need repeating.

This failure is the scientific opening for the present work. If a different
architecture or a different input representation fixes the random-translation
exponent, that is the result. It is a target, not a bug to be papered over.
% <<< END BLOCK: failure-mode

% =====================================================================
\section{Reproduced results (living record)}
\label{sec:reproduced}
% =====================================================================

% >>> BLOCK: table-phase1
\subsection{Phase 1 --- classifiers in isolation}

Classification accuracy on the two \emph{extreme} regimes only. These numbers
establish that the pipeline is sound; they are not a physics result, because
the two classes are trivially separable at the extremes.

\begin{table}[h]
\centering
\small
\begin{tabular}{llccccc}
\toprule
\textbf{Model} & \textbf{Channel} & $N$ & \textbf{Samples} & \textbf{Train} & \textbf{Test} & \textbf{Provenance} \\
\midrule
SVM & discrete    & 80 & 1800 & --- & $1.0000$ & \texttt{configs/svm\_discrete\_coin.json} \\
SVM & continuous  & 80 & 1800 & --- & $1.0000$ & \texttt{configs/svm\_continuous\_coin.json} \\
SVM & random tr.  & 80 & 1800 & --- & $1.0000$ & \texttt{configs/svm\_random\_translation.json} \\
MLP & discrete    & 80 & 1800 & $1.0000$ & $1.0000$ & regime \texttt{discrete\_coin}, seed 0 \\
MLP & continuous  & 80 & 1800 & $1.0000$ & $1.0000$ & regime \texttt{continuous\_coin}, seed 0 \\
MLP & random tr.  & 80 & 1800 & $1.0000$ & $1.0000$ & regime \texttt{random\_translation}, seed 0 \\
MLP & discrete    & 300 & 600 & $1.0000$ & $1.0000$ & regime \texttt{large\_lattice}, seed 0 \\
CNN & ---         & --- & --- & ---      & ---      & not yet implemented \\
\bottomrule
\end{tabular}
\caption{Phase-1 accuracies. All MLP rows: architecture $(400,200,100,50)$,
$\alpha=0.001$, $P_{\max}=1$ normalization on, \texttt{random\_state}~$=0$,
$80/20$ split, produced by \texttt{scripts/train\_mlp.py}. Windows are the
placeholders recorded in that script's \texttt{REGIMES} table and are
\emph{not} yet tuned.}
\label{tab:phase1}
\end{table}
% <<< END BLOCK: table-phase1

% >>> BLOCK: table-exponents
\subsection{Phases 2--3 --- critical values and exponents}

Empty until Phase 2 produces confusion points. Keep the column structure of
Table~\ref{tab:paperA} so the two can be set side by side.

\begin{table}[h]
\centering
\begin{tabular}{lccc}
\toprule
\textbf{Method} & $\dthetac$ (discrete) & $\dthetac$ (continuous) & $\Prandc$ \\
\midrule
SVM    & --- & --- & --- \\
MLP NN & --- & --- & --- \\
CNN    & --- & --- & --- \\
MoI    & --- & --- & --- \\
Human  & --- & --- & --- \\
IPR    & --- & --- & --- \\
\bottomrule
\end{tabular}
\caption{This work. Exponents $\alpha$ in $\dthetac \sim N^{-\alpha}$, to be
compared against Table~\ref{tab:paperA}. Each entry must carry the fit range
in $N$, the training windows, and the seed.}
\label{tab:exponents}
\end{table}
% <<< END BLOCK: table-exponents

% =====================================================================
\section{Observations and open questions}
\label{sec:observations}
% =====================================================================

% >>> BLOCK: observations
\begin{enumerate}

\item \textbf{Ballistic versus diffusive spreading, measured.}
Fitting $\MoI \sim N^{\alpha}$ over $N \in [8,300]$ gives
$\alpha_{\mathrm{quantum}} = 1.996$ and $\alpha_{\mathrm{classical}} = 1.000$,
against the expected $2$ and $1$. At $N=300$, $\theta=\pi/4$, the quantum
walk gives $\MoI/N^{2} = 0.29290$, which is $1-1/\sqrt{2} = 0.29289$ to five
digits --- the standard Hadamard-walk variance coefficient. This is a
correctness check on the evolution code as much as a physics statement.
\emph{Provenance:} \texttt{scripts/legacy\_classical\_walk.py}, defaults,
seed 42.

\item \textbf{Phase-1 accuracy saturates, so Phase-1 grid search is
uninformative.}
Every point of the $3\times 3$ architecture $\times$ $\alpha$ grid scores CV
accuracy $1.0000$ with standard deviation $0$ on the \texttt{discrete\_coin}
regime. Grid search cannot rank hyperparameters on a problem every
hyperparameter solves. Tuning is only meaningful once there is a
transition-regime score to select on --- which is Phase 2. This is a
sequencing consequence, not a defect.

\item \textbf{Where the two-class problem stops being trivial.}
Squeezing the two training windows toward each other on
\texttt{discrete\_coin} at $N=80$, test accuracy stays exactly $1.0000$ until
the gap closes to roughly $\dtheta \in [0.22, 0.28]$, then falls
($0.9917$ at a $[0,0.22]$ vs $[0.28,\theta_0]$ split, $0.9667$ at
$[0,0.24]$ vs $[0.26,\theta_0]$).
\provisional{This brackets the region the classifier cannot separate, near
$\dtheta \approx 0.23$--$0.25$ at $N=80$. It is a Phase-1 diagnostic and is
\emph{not} a measurement of $\dthetac$ --- the confusion-point rules of
Sec.~\ref{sec:critical} have not been applied. Do not quote it as a critical
value.}
\emph{Provenance:} \texttt{scripts/train\_mlp.py}, 600 samples, seed 0,
\texttt{random\_state}~$=0$.

\item \textbf{On the squeezed windows, Paper~A's architecture does come out
on top --- but not significantly.}
Using the deliberately hard \texttt{stress} regime, the coarse grid ranks
$(400,200,100,50)$ first at CV accuracy $0.9583 \pm 0.0106$, against
$0.9542 \pm 0.0118$ for $(100,50)$. The error bars overlap almost completely,
so this is consistent with Paper~A's choice but is not independent evidence
for it. More folds or more seeds would be needed to say anything stronger.

\item \textbf{The MLP's first-layer sensitivity respects the parity comb.}
The per-site $\ell_2$ norm of the first-layer weights is near zero on the
unreachable (identically zero) sites and largest near the ballistic edge,
$|x| \approx 65$ at $N=80$. The network is keying on the edge of the light
cone, which is the physically sensible feature. Note this is a
\emph{sensitivity}, not a decision boundary --- the MLP is nonlinear, so
unlike the linear SVM's weights it cannot be read directionally.

\item \textbf{Open --- random translation is suspiciously easy so far.}
The \texttt{random\_translation} regime scores $1.0000$ like everything else.
At the extremes ($\Prand \approx 0$ versus $\Prand \approx 0.5$) that is
expected and benign. The known failure mode is a transition-regime effect,
driven by the three-peak structure of Sec.~\ref{sec:phenomenology}, so it
cannot appear until Phase 2. Flagged here so that a clean Phase-1 number is
not mistaken for the problem being absent.

\end{enumerate}
% <<< END BLOCK: observations

% =====================================================================
\section{Figure inventory}
\label{sec:figures}
% =====================================================================

% >>> BLOCK: figures
\begin{longtable}{p{0.40\textwidth} p{0.54\textwidth}}
\toprule
\textbf{File in \texttt{figures/}} & \textbf{Content} \\
\midrule
\endhead
\texttt{Total\_probabilities.png}
  & Hadamard QW, 300 steps, total $P(x)$; clean two-peak. \\
\texttt{Spin\_probabilities.png}
  & $P_+$ (blue solid) against $P_-$ (red dashed), 300 steps; peaks near
    $x \approx \pm 210$, components avoid each other. \\
\texttt{Pure\_Coin\_vs\_\_Random\_Coin.png}
  & Pure against discrete-random and continuous-random coins, 300 steps. \\
\texttt{Randoms\_vs\_\_Control.png}
  & The above plus random translation, 300 steps. Defines the project's line
    styles. \\
\texttt{Peak\_spin\_updown\_theta\_vs\_distance.png}
  & Peak spin-up position against coin angle $\theta \in [-\pi/2,\pi/2]$;
    maximum $\approx 300$ at $\theta=0$, falling to $0$ at $\theta=\pm\pi/2$. \\
\texttt{classical\_walk\_distribution.png}
  & Classical binomial against pure QW at 300 steps, drawn on the reachable
    sub-lattice, with an optional Monte Carlo overlay. \\
\texttt{classical\_walk\_spreading.png}
  & $\MoI$ against $N$, log--log; fitted slopes $1.996$ (quantum) and
    $1.000$ (classical). \\
\texttt{mlp\_stress\_coarse\_grid.png}
  & MLP diagnostic panel: training classes, loss curve, first-layer
    sensitivity, coarse-grid CV scores. Diagnostic regime --- not quotable. \\
\bottomrule
\end{longtable}
% <<< END BLOCK: figures

% =====================================================================
\section*{Source papers}
\addcontentsline{toc}{section}{Source papers}
% =====================================================================

\begin{thebibliography}{9}

\bibitem{PaperA}
C.~Mastandrea and C.-C.~Chien,
\emph{Localization of quantum walks with classical randomness: Comparison
between manual methods and supervised machine learning},
Phys.\ Rev.\ E \textbf{108}, 035308 (2023).

\bibitem{PaperB}
C.~Mastandrea and C.-C.~Chien,
\emph{Robustness and classical proxy of entanglement in variants of quantum
walks},
Phys.\ Rev.\ E \textbf{110}, 064124 (2024).

\end{thebibliography}

\end{document}
